\clearpage\corrigeDevoir{1}
\begin{blocLivre}{dscor-1-1}\begin{corrigeBac}{6}{Identifier un liquide}
\demarche{Soustraire la tare, puis comparer des masses volumiques dans les mêmes conditions.}
\begin{questions}
\item $m=54,9-35,2=19,7$ g ; $\rho=19,7/25,0=\qty{0,788}{g.mL^{-1}}$.\pts{2 pt}
\item L’éthanol est compatible ; une masse volumique ne suffit pas à exclure un mélange ou une autre espèce proche.\pts{2 pt}
\item La moyenne vaut \qty{0,79}{g.mL^{-1}}. Employer une verrerie mieux adaptée, lire le ménisque à hauteur des yeux et contrôler la température.\pts{2 pt}
\end{questions}
\end{corrigeBac}\end{blocLivre}
\begin{blocLivre}{dscor-1-2}\begin{corrigeBac}{7}{Une préparation précise}
\demarche{Conserver la distinction entre masse de soluté et volume final de solution.}
\schema[corrige]{dissolution}
\begin{questions}
\item L’eau est le solvant.\pts{1 pt}
\item $m=c_mV=12,0\times0,2000=\qty{2,40}{g}$.\pts{2 pt}
\item Peser, dissoudre, transférer et rincer dans une fiole de 200,0 mL, compléter au trait et homogénéiser.\pts{2 pt}
\item Le volume est trop grand, donc la concentration trop faible. Retirer un mélange homogène retire solvant et soluté dans la même proportion, sans modifier la concentration.\pts{2 pt}
\end{questions}
\end{corrigeBac}\end{blocLivre}
\begin{blocLivre}{dscor-1-3}\begin{corrigeBac}{7}{Diluer avant de doser}
\demarche{Lire la solution effectivement mesurée, puis remonter à la solution initiale.}
\schema[corrige]{etalonnage}
\begin{questions}
\item $c_f=\qty{60}{g.L^{-1}}$ par projection sur la droite.\pts{2 pt}
\item $F=100,0/20,0=5,00$ ; $c_m=5,00\times60=\qty{300}{g.L^{-1}}$ dans la précision de lecture.\pts{2 pt}
\item Une pipette jaugée de 20,0 mL avec propipette.\pts{1 pt}
\item Le modèle d’étalonnage n’est validé que sur la gamme ; il faut diluer davantage ou préparer une gamme adaptée.\pts{2 pt}
\end{questions}
\end{corrigeBac}\end{blocLivre}

